% 第七章：一致收敛性

\section{一致收敛性}

数二的级数部分涉及逐点收敛，即对每个固定的 $x$，级数 $\sum_{n=1}^\infty u_n(x)$ 是否收敛。数一额外引入了一致收敛的概念，它要求级数不仅在每一点收敛，而且收敛的"速度"在整个定义域上是均匀的。一致收敛性是函数项级数理论中最核心的概念，它决定了极限运算（求和、求导、积分）是否可以交换顺序——这是数学分析中最深刻的问题之一。

\subsection{逐点收敛与一致收敛的区别}

\subsubsection{逐点收敛}

函数项级数 $\sum_{n=1}^\infty u_n(x)$ 在区间 $I$ 上\textbf{逐点收敛}，是指对每个 $x\in I$，数项级数 $\sum_{n=1}^\infty u_n(x)$ 收敛。记部分和为 $S_n(x) = \sum_{k=1}^n u_k(x)$，和函数为 $S(x)$，则逐点收敛意味着：
\[
\forall x\in I,\ \forall\varepsilon>0,\ \exists N=N(x,\varepsilon),\ \text{s.t.}\ n>N \Rightarrow |S_n(x)-S(x)|<\varepsilon
\]

注意：$N$ 依赖于 $x$ 和 $\varepsilon$。不同的 $x$ 可能需要不同的 $N$ 才能使 $|S_n(x)-S(x)|<\varepsilon$。

\subsubsection{一致收敛}

函数项级数 $\sum_{n=1}^\infty u_n(x)$ 在区间 $I$ 上\textbf{一致收敛}，是指：
\[
\forall\varepsilon>0,\ \exists N=N(\varepsilon),\ \text{s.t.}\ n>N \Rightarrow |S_n(x)-S(x)|<\varepsilon,\ \forall x\in I
\]

关键区别：$N$ 只依赖于 $\varepsilon$，不依赖于 $x$。也就是说，存在一个"统一的"$N$，使得对所有 $x$，$n>N$ 时 $|S_n(x)-S(x)|<\varepsilon$ 同时成立。

\begin{essencebox}
一致收敛的本质是\textbf{收敛的均匀性}。逐点收敛只保证在每个点上级数收敛，但不同点的收敛速度可能差异巨大——有些点很快收敛，有些点很慢。一致收敛要求收敛速度在整个区间上"均匀"——不存在收敛特别慢的"拖后腿"的点。

用 $\varepsilon$-带的语言来说：一致收敛意味着，对任意 $\varepsilon>0$，存在 $N$，使得 $n>N$ 时，$S_n(x)$ 的图像完全落在 $S(x)\pm\varepsilon$ 的带状区域内。而逐点收敛只保证对每个 $x$，$S_n(x)$ 最终进入带状区域，但进入的"时刻"可能因 $x$ 而异。
\end{essencebox}

\begin{figure}[htbp]
\centering
\begin{tikzpicture}
\begin{axis}[
    width=10cm, height=7cm,
    xlabel={$x$}, ylabel={$y$},
    domain=0.01:1, samples=100,
    legend style={at={(0.02,0.98)}, anchor=north west, font=\small},
    ymin=-0.1, ymax=1.1
]
% S(x) = x
\addplot[blue, very thick] {x};
\addlegendentry{$S(x)=x$}
% epsilon band
\addplot[gray, dashed, thin] {x+0.1};
\addplot[gray, dashed, thin] {x-0.1};
\addlegendentry{$\varepsilon$-带}
% S_n(x) = x^n (pointwise but not uniform)
\addplot[red, thick] {x^5};
\addlegendentry{$S_5(x)=x^5$}
\addplot[orange, thick] {x^10};
\addlegendentry{$S_{10}(x)=x^{10}$}
\end{axis}
\end{tikzpicture}
\caption{逐点收敛但不一致收敛的例子：$S_n(x)=x^n$ 在 $[0,1]$ 上逐点收敛到 $0$（$x<1$）和 $1$（$x=1$），但收敛不一致}
\end{figure}

\subsection{一致收敛的判别法}

\subsubsection{Weierstrass 判别法（优级数法）}

若存在正项级数 $\sum M_n$ 使得 $|u_n(x)| \leq M_n$ 对所有 $x\in I$ 成立，且 $\sum M_n$ 收敛，则 $\sum u_n(x)$ 在 $I$ 上一致收敛。

\begin{insightbox}
\textbf{推理步骤}：Weierstrass 判别法的逻辑
\begin{enumerate}[leftmargin=2em]
    \item $|S_n(x)-S(x)| = |\sum_{k=n+1}^\infty u_k(x)| \leq \sum_{k=n+1}^\infty |u_k(x)| \leq \sum_{k=n+1}^\infty M_k$
    \item 因为 $\sum M_n$ 收敛，所以 $\sum_{k=n+1}^\infty M_k \to 0$（当 $n\to\infty$）
    \item 因此 $|S_n(x)-S(x)| \leq \sum_{k=n+1}^\infty M_k$ 对所有 $x$ 成立
    \item 取 $N$ 使得 $n>N$ 时 $\sum_{k=n+1}^\infty M_k < \varepsilon$，则 $|S_n(x)-S(x)|<\varepsilon$ 对所有 $x$ 成立
    \item 这正是一致收敛的定义
\end{enumerate}
Weierstrass 判别法的核心思想是：用一个收敛的数项级数来"控制"函数项级数，从而保证收敛的均匀性。
\end{insightbox}

\subsection{一致收敛的性质}

\subsubsection{连续性传递}

若 $u_n(x)$ 在 $[a,b]$ 上连续，且 $\sum u_n(x)$ 在 $[a,b]$ 上一致收敛，则和函数 $S(x)$ 在 $[a,b]$ 上连续。

\subsubsection{逐项积分}

若 $u_n(x)$ 在 $[a,b]$ 上连续，且 $\sum u_n(x)$ 在 $[a,b]$ 上一致收敛，则：
\[
\int_a^b \sum_{n=1}^\infty u_n(x)\,\mathrm{d}x = \sum_{n=1}^\infty \int_a^b u_n(x)\,\mathrm{d}x
\]

即\textbf{积分与求和可以交换顺序}。

\subsubsection{逐项求导}

若 $u_n(x)$ 在 $[a,b]$ 上有连续导数，$\sum u_n'(x)$ 在 $[a,b]$ 上一致收敛，且 $\sum u_n(x)$ 至少在某一点收敛，则 $\sum u_n(x)$ 在 $[a,b]$ 上一致收敛，且：
\[
\frac{\mathrm{d}}{\mathrm{d}x}\sum_{n=1}^\infty u_n(x) = \sum_{n=1}^\infty u_n'(x)
\]

即\textbf{求导与求和可以交换顺序}。

\begin{essencebox}
逐项求导的条件比逐项积分更强——不仅要求 $\sum u_n'(x)$ 一致收敛，还要求 $\sum u_n(x)$ 至少在某一点收敛。这是因为求导是"放大差异"的运算（微分使函数变得更粗糙），而积分是"平滑差异"的运算（积分使函数变得更光滑）。因此，逐项求导需要更强的条件来保证合理性。

这三种交换性——连续性传递、逐项积分、逐项求导——是一致收敛性最核心的价值。它们回答了数学分析中的一个根本问题：\textbf{什么时候可以交换极限运算的顺序？} 这个问题贯穿了整个分析学，从级数理论到勒贝格积分，从傅里叶分析到分布理论。
\end{essencebox}

\subsection{一致收敛与幂级数}

幂级数 $\sum_{n=0}^\infty a_n(x-x_0)^n$ 在其收敛区间内的任何闭子区间上一致收敛（由 Weierstrass 判别法可证）。这意味着：

\begin{enumerate}[leftmargin=2em]
    \item 幂级数的和函数在收敛区间内连续
    \item 幂级数在收敛区间内可以逐项积分
    \item 幂级数在收敛区间内可以逐项求导，且求导后的幂级数收敛半径不变
\end{enumerate}

\begin{connectionbox}
\textbf{1. 统一极限运算的交换理论}

一致收敛性为"交换极限顺序"提供了充分条件。在数二中，我们不加证明地使用了"幂级数可以逐项求导和逐项积分"的结论。理解了一致收敛性后，这个结论不再是"记住就好"的规则，而是有严格证明的定理。这种从"知其然"到"知其所以然"的转变，是数学理解深化的标志。

\textbf{2. 深化对收敛概念的理解}

逐点收敛和一致收敛的区分，揭示了"收敛"这一概念的层次性。逐点收敛是"最弱"的收敛，只保证在每个点上级数收敛；一致收敛是"更强"的收敛，保证收敛的均匀性。在更高级的分析学中，还有更强的收敛概念（如 $L^p$ 收敛、几乎处处收敛等），它们构成了收敛概念的谱系。

\textbf{3. 为傅里叶级数提供理论支撑}

傅里叶级数的逐项积分不需要一致收敛的条件——这是傅里叶级数的一个特殊性质。但傅里叶级数的逐项求导需要一致收敛（或更强的条件）。理解了一致收敛性，就能理解为什么傅里叶级数可以逐项积分但不能随意逐项求导。

\textbf{4. 促进对连续与可微关系的理解}

一致收敛保证和函数连续，但不保证和函数可微。要保证和函数可微，需要对导数级数加一致收敛的条件。这深化了我们对"连续 $\not\Rightarrow$ 可微"这一基本事实的理解——即使每个 $u_n$ 都光滑，极限函数也可能不光滑，除非收敛的"质量"足够好（一致收敛）。

\textbf{5. 为函数空间理论奠定基础}

一致收敛性本质上是在函数空间 $C[a,b]$（连续函数空间）中定义的一种拓扑。在这个空间中，一致收敛就是"按上确界范数收敛"。这种观点将级数理论与泛函分析联系起来，为理解 Banach 空间、完备性等高级概念提供了直觉基础。
\end{connectionbox}
